% Section 1.1, from lectures/L01/appendix/a_c_fundamentals.md.

\section{C Fundamentals for Experienced Programmers}
\label{c1:sec:fundamentals}\label{c1:app:a}
A fast-moving refresher on C syntax for programmers who already know how to program in another
language; not a full C course, but enough that later chapters can assume you're not fighting the
language while learning hardware concepts. If a section looks obvious, skip ahead; the pointer
section (\secref{c1:a:7}) is the one worth actually reading.

\subsection{Compiling and program shape}
\label{c1:a:1}
\begin{itemize}
  \item C is compiled, not interpreted: source code $\rightarrow$ assembly $\rightarrow$ machine
    code. No runtime interpretation overhead, which is most of why it's the default choice for
    embedded systems.
  \item C is procedural: everything lives in functions; \code{main()} is the entry point.
  \item No safety net: no bounds checking, no garbage collection. The language trusts you.
\end{itemize}
A \code{Hello, world!} program is shown below:

\begin{ccode}
#include <stdio.h>

int main(void)
{
    printf("Hello, world!\n");
    return 0;
}
\end{ccode}

\begin{itemize}
  \item \code{#include <stdio.h>} pulls in declarations (a header), not an implementation; that's
    resolved at link time.
  \item \code{return 0} from \code{main()} signals success to the OS. C99+ makes it implicit if
    omitted, but write it anyway.
  \item A Doxygen comment above a function, written in the imperative mood and documenting its
    contract, is conventional throughout this codebase. For example, a function that computes
    voltage by Ohm's law can be documented as shown below:
\end{itemize}

\begin{ccode}
/**
 * @brief Compute voltage from resistance and current.
 *
 * @param[in] resistance_kohm Resistance in kOhm.
 * @param[in] current_ma Current in mA.
 *
 * @return Voltage in V.
 */
double compute_voltage(double resistance_kohm, double current_ma);
\end{ccode}

\subsection{Data types \& naming}
\label{c1:a:2}

\begin{booktable}{@{}p{9.6em}Lp{6.4em}@{}}
  \thead{Type} & \thead{Meaning} & \thead{Typical size} \\
  \midrule
  \code{void} & nothing (no return value) & -- \\
  \code{int} & signed integer & 2--4 bytes \\
  \code{float} & single-precision float & 4 bytes \\
  \code{double} & double-precision float (default for float literals) & 8 bytes (4 with avr-gcc) \\
  \code{char} & one character & 1 byte \\
  \code{char[]} / \code{char*} & string (array of \code{char}, NUL-terminated) & -- \\
  \code{bool} (\code{<stdbool.h>}) & \code{true} / \code{false} & 1 byte \\
\end{booktable}

\begin{itemize}
  \item Prefer \code{const} for anything that shouldn't change after initialization, following the
    same intent as \code{final}/\code{readonly} in other languages. It also lets the compiler catch
    mistakes and optimize more aggressively.
\end{itemize}
An exception is function declarations: do not qualify arguments passed by value with \code{const},
as it has no effect on the function interface and only reduces readability. In the corresponding
function definition, however, pass-by-value arguments should be declared \code{const} for the
reasons above.
\begin{itemize}
  \item Naming convention: \code{snake_case} for variables/functions/constants, \code{UPPER_CASE}
    for macros and enum constants. Not camelCase/PascalCase; that's a Java/C\#/Pascal/C++
    convention, not a C one.
  \item \textbf{Integer division truncates}: \code{5 / 2 == 2}, not \code{2.5}. Force float division
    by making at least one operand a \code{double}/\code{float}, or cast: \code{(double)(a) / b}.
  \item \code{\%} is modulus, integers only; using it on a float is a compiler error.
\end{itemize}

\subsection{Structs, enums, and typedefs}
\label{c1:a:3}

\subsubsection{Structs}
\begin{itemize}
  \item A \code{struct} groups related fields under one name, each accessed with \code{.} (or
    \code{->} if you're going through a pointer). No methods, no inheritance, no access control;
    just a memory layout, laid out in declaration order (the compiler may insert padding between
    fields for alignment).
  \item Compare to a plain data class/record in other languages, minus the behavior.
  \item The idiom used throughout this codebase for functions that operate on a struct: take a
    pointer to it as the first parameter, named \code{self}. It's the closest C gets to a method's
    implicit \code{this}, without the language actually having objects:
\end{itemize}

\begin{ccode}
/**
 * @brief Electrical measurement structure.
 */
struct measurement
{
    /** Voltage in V. */
    double voltage_v;

    /** Current in mA. */
    double current_ma;

    /** Resistance in kOhm. */
    double resistance_kohm;
};

/**
 * @brief Store a voltage/current reading and derive its resistance (Ohm's law).
 *
 * @param[out] self Pointer to the measurement to populate.
 * @param[in] voltage_v Voltage in V.
 * @param[in] current_ma Current in mA.
 */
static void measurement_init(struct measurement* self, const double voltage_v,
                              const double current_ma)
{
    self->voltage_v       = voltage_v;
    self->current_ma      = current_ma;
    self->resistance_kohm = voltage_v / current_ma;
}

struct measurement m;
measurement_init(&m, 5.0, 0.5);
printf("%g V\n", m.voltage_v);
\end{ccode}

\begin{itemize}
  \item Passing a \code{struct} by value copies every field, so \code{self} is (almost) always a
    pointer:
  \begin{itemize}
    \item A \code{const struct measurement*} when the function only reads it.
    \item A plain \code{struct measurement*} when it writes through it, same as any other
      out-parameter (\secref{c1:a:7}).
  \end{itemize}
\end{itemize}

\subsubsection{\code{typedef}}
\begin{itemize}
  \item \code{typedef} gives an existing type a new name, most often to avoid typing
    \code{struct measurement} everywhere. The convention in this codebase is a \code{_t} suffix on
    an anonymous struct:
\end{itemize}

\begin{ccode}
/**
 * @brief Electrical measurement structure.
 */
typedef struct
{
    /** Voltage in V. */
    double voltage_v;

    /** Current in mA. */
    double current_ma;

    /** Resistance in kOhm. */
    double resistance_kohm;
} measurement_t;

// Create measurement instance.
measurement_t m = { .voltage_v = 5.0, .current_ma = 0.5, .resistance_kohm = 10.0 };
\end{ccode}

\begin{itemize}
  \item \code{measurement_t} behaves exactly like \code{struct measurement} above; the
    \code{typedef} is purely a naming convenience, it doesn't change layout or behavior.
\end{itemize}

\subsubsection{\code{enum}}
\begin{itemize}
  \item An \code{enum} names a set of related integer constants, so a \code{switch} reads as intent
    rather than magic numbers, and some compilers can warn you when a \code{switch} doesn't handle
    every named case:
\end{itemize}

\begin{ccode}
/**
 * @brief Enumeration of sensor states.
 */
typedef enum
{
    SENSOR_STATE_IDLE,     ///< Sensor is idle.
    SENSOR_STATE_SAMPLING, ///< Sensor is sampling.
    SENSOR_STATE_ERROR     ///< Sensor is in an error state.
} sensor_state_t;

// Create variable holding the sensor state.
sensor_state_t state = SENSOR_STATE_IDLE;
\end{ccode}

\begin{itemize}
  \item Unless you assign values explicitly, members count up from 0 in declaration order
    (\code{SENSOR_STATE_IDLE == 0}, \code{SENSOR_STATE_SAMPLING == 1}, \dots).
  \item Assign explicit values when the numbers themselves matter, e.g.\ they're written to a
    hardware register.
  \item Prefer this over a block of \code{#define SENSOR_STATE_IDLE 0} macros: same cost at
    runtime, but the type shows up in a debugger and the compiler can help catch an unhandled
    \code{switch} case.
\end{itemize}

\subsection{Scope, lifetime, and storage}
\label{c1:a:4}

\begin{booktable}{@{}p{7.4em}LL@{}}
  \thead{Storage class} & \thead{Scope} & \thead{Lifetime} \\
  \midrule
  \textbf{Local} & Inside a function or block. & Stack-allocated; destroyed when the block
    exits. \\
  \textbf{\code{static} local} & Inside a function or block. & Persists for the program's life. \\
  \textbf{Global} & Outside any function. & Exists for the program's life; stored in RAM. \\
\end{booktable}

\noindent\textbf{Rule of thumb}:
\begin{itemize}
  \item Minimize scope and lifetime.
  \item Mark file-scope functions and globals \code{static} when they shouldn't leak outside the
    \code{.c} file; it's the closest thing C has to ``private,'' and it gives the compiler more
    room to inline and optimize.
\end{itemize}

\subsection{Printing (\code{printf()})}
\label{c1:a:5}

\begin{booktable}{@{}p{7em}L@{}}
  \thead{Specifier} & \thead{Type} \\
  \midrule
  \code{\%d} / \code{\%i} & \code{int} \\
  \code{\%u} & unsigned int \\
  \code{\%f} & \code{double} (a \code{float} argument is promoted to \code{double}), fixed
    decimals; \code{\%lf} means the same \\
  \code{\%g} & \code{double}, shortest representation \\
  \code{\%c} & single \code{char} \\
  \code{\%s} & string (\code{char*}) \\
  \code{\%p} & pointer address \\
\end{booktable}

\begin{ccode}
printf("%s is %d years old!\n", name, age);
printf("Voltage: %.2f V\n", voltage);   // 2 decimal places
\end{ccode}

\subsection{Functions}
\label{c1:a:6}

\begin{ccode}
return_type function_name(param_type param_name, ...)
{
    ...
    return value;   // omit for void
}
\end{ccode}

\begin{itemize}
  \item \code{static} before the return type = internal linkage (only visible in this \code{.c}
    file). Default to \code{static} for anything that isn't part of a module's public API.
  \item \code{static inline} on small, frequently-called functions hints the compiler to paste the
    body in at the call site instead of paying for a call/return. Modern compilers usually do this
    anyway, but it's idiomatic for small helpers.
  \item Declare functions (a prototype) before use, typically in a header, so callers don't need
    the full definition.
\end{itemize}

\begin{ccode}
/**
 * @brief Compute resistance from voltage and current (Ohm's law).
 *
 * @param[in] voltage Voltage in V.
 * @param[in] current Current in mA.
 *
 * @return Resistance in kOhm, or 0 if the current is 0.
 */
static inline double get_resistance(const double voltage, const double current)
{
    return 0.0 != current ? voltage / current : 0.0;
}
\end{ccode}

\begin{itemize}
  \item Convention for ordering a \code{.c} file: forward-declare its \code{static} helpers once,
    right after the \code{#include}s, then define the public functions (in the same order the
    header declares them), and put the \code{static} helpers' own definitions at the bottom of the
    file. A reader hits the public API first and only descends into private implementation detail
    if they keep scrolling:
\end{itemize}

\begin{ccode}
#include "sensor.h"

/* Static functions: */
static double filter_reading(double raw_v);

double sensor_read(void)
{
    const double raw_v = adc_read();
    return filter_reading(raw_v);
}

/* raw_v filtering is an implementation detail, not part of sensor.h. */
static double filter_reading(const double raw_v)
{
    ...
}
\end{ccode}

\begin{itemize}
  \item When a header can be included from more than one place (or pulled in indirectly through
    another header), wrap its contents in a header guard so the compiler never sees the same
    declarations twice:
\end{itemize}

\begin{ccode}
#ifndef SENSOR_H_
#define SENSOR_H_

double sensor_read(void);

#endif /* SENSOR_H_ */
\end{ccode}

\begin{itemize}
  \item Convention: the guard macro is the file name in \code{UPPER_CASE}, with \code{.} replaced
    by \code{_} and a trailing \code{_} (\file{sensor.h} $\rightarrow$ \code{SENSOR_H_}), chosen so
    it's unlikely to collide with any other macro in the program. \code{#pragma once} does the same
    job in a single line and most modern compilers support it, but it isn't part of the C standard;
    traditional guards are the portable choice and what you'll see throughout this codebase when
    using C. In C++, \code{#pragma once} is also widely used.
\end{itemize}

\subsection{Pointers (the one thing that's actually different)}
\label{c1:a:7}

\begin{aside}
\textbf{Everything is pass-by-value in C.} Every function argument is copied. Passing a pointer
still passes a copy of the pointer; it just happens to point to the same object as the caller's
pointer. This lets a function modify the caller's object through the pointer, but assigning a new
value to the pointer itself does not affect the caller.
\end{aside}

\noindent If you're coming from a garbage-collected language, this is the section to actually read:
\begin{itemize}
  \item A pointer holds an \emph{address}, e.g.\ \code{int* ptr = &x;} makes \code{ptr} point at
    \code{x}:
  \begin{itemize}
    \item \code{&x} means ``address of x''.
    \item \code{*ptr} means ``the value that \code{ptr} points to'' (dereference).
  \end{itemize}
  \item Pointers let a function modify the caller's object because, although the pointer itself is
    passed by value, it still points to that object.
\end{itemize}

\begin{ccode}
/**
 * @brief Assign 3 to the variable pointed to by x.
 *
 * @param[out] x Pointer to the variable to assign.
 */
static void assign(int* x)
{
    *x = 3; // Writes through the pointer into the caller's variable.
}

int main(void)
{
    int a = 0;
    assign(&a);
    printf("a = %d\n", a); // a = 3.
}
\end{ccode}

\noindent This is the standard C idiom for ``multiple out-parameters'' or ``modify in place,''
using no references, no boxing, just an address. You'll see this pattern constantly once we start
writing to hardware registers:
\begin{itemize}
  \item A pointer to a pointer (\code{int**}) lets a function replace \emph{what the caller's
    pointer points to}, not just the value at that address; the same out-parameter idiom, one level
    deeper. You'll see this whenever a function needs to null out or reassign a pointer it was only
    given a copy of:
\end{itemize}

\begin{ccode}
#include <stddef.h>

static void set_to_null(int** ptr)
{
    *ptr = NULL; // Changes the caller's pointer itself, not just *ptr's target.
}

int main(void)
{
    int x = 5;
    int* p = &x;
    set_to_null(&p);
    // p is now NULL; x is untouched.
}
\end{ccode}

\begin{itemize}
  \item Arrays and strings decay to a pointer to their first element automatically when passed to
    a function, no \code{&} needed:
\end{itemize}

\begin{ccode}
static void print_text(const char* s) { printf("%s", s); }
const char text[] = "no & needed here";
print_text(text);
\end{ccode}

\begin{itemize}
  \item \code{const char*} means ``pointer to data I won't modify through this pointer''; use it on
    any string/array parameter you're only reading, same spirit as marking a parameter
    \code{readonly}.
  \item A string is a \code{char} array ending in a \code{'\0'} (NUL) byte; there's no length
    field. \code{sizeof} gives you the buffer's declared capacity, not the content length, and only
    works on the actual array, not a pointer to it.
\end{itemize}

\subsection{Function pointers}
\label{c1:a:8}
A function pointer holds the \emph{address of a function} instead of a variable, callable through
the pointer exactly like calling the function by name:

\begin{ccode}
#include <stdio.h>

static void greet1(void) { printf("Hello!\n"); }
static void greet2(void) { printf("World!\n"); }

int main(void)
{
    void (*fn)(void) = greet1; // fn now holds greet1's address.
    fn();                      // Calls greet1(), exactly like greet1() would.

    fn = greet2;               // fn now holds greet2's address.
    fn();                      // Calls greet2(), exactly like greet2() would.
}
\end{ccode}

\begin{itemize}
  \item The declaration reads inside-out: \code{void (*fn)(void)} is ``\code{fn} is a pointer to a
    function taking no arguments and returning \code{void}.'' The parentheses around \code{*fn} are
    required, \code{void *fn(void)} would instead declare a function returning \code{void*}.
  \item A function's name, used without calling it (no \code{()}), decays to its address the same
    way an array name decays to a pointer to its first element (\secref{c1:a:7}):
    \code{fn = greet1;} needs no \code{&}.
  \item \code{typedef} cleans up the syntax the same way it does for structs (\secref{c1:a:3}), and
    gives the pointer type a name that documents what it's for:
\end{itemize}

\begin{ccode}
typedef void (*callback_t)(void);

static void on_timeout(void) { printf("Timed out!\n"); }

static void run(callback_t cb) // cb holds a function's address, not a value.
{
    cb(); // Calls whichever function was passed in.
}

int main(void)
{
    run(on_timeout); // Prints "Timed out!".
}
\end{ccode}

\begin{itemize}
  \item \code{run()} never mentions \code{on_timeout} by name, it just calls whatever
    \code{callback_t} it was given. That's the whole idea: the caller decides which function runs,
    the callee just knows the function's signature.
  \item This comes back twice, later in the book: Chapter~9's timer driver exercise
    (Exercise~\ref{c9:ex:6}) stores one of these in a struct so a timer can run arbitrary code when
    it expires, and Chapter~\ref{c10:ch} builds an entire polymorphism system, structs whose
    behavior is decided at runtime rather than compile time, out of nothing but function pointers.
\end{itemize}

\subsection{Reading input}
\label{c1:a:9}
You won't have a keyboard on a microcontroller, but the buffer-safety idiom here reappears when
reading serial/UART data later, so it's worth knowing.
\begin{itemize}
  \item \code{scanf("\%s", buf)} is easy but unsafe: no bounds checking, and it stops at the first
    whitespace.
  \item \code{fgets(buf, sizeof(buf), stdin)} is the safe alternative; it takes an explicit buffer
    size and won't overflow, but it keeps the trailing \code{'\n'}, which callers usually strip:
\end{itemize}

\begin{ccode}
/**
 * @brief Read a line of input into a buffer, stripping the trailing newline.
 *
 * @param[out] buf Buffer to store the line in.
 * @param[in] buflen Capacity of the buffer, in bytes.
 */
static void readline(char* buf, const size_t buflen)
{
    // Read data from the terminal, store in the buffer.
    fgets(buf, (int)(buflen), stdin);

    // Replace '\n' characters with '\0'.
    for (size_t i = 0U; '\0' != buf[i]; i++)
    {
        if ('\n' == buf[i]) { buf[i] = '\0'; }
    }
}
\end{ccode}

\begin{itemize}
  \item To get numbers out of text: read a line as a string, then convert with \code{atoi()} for
    \code{int} or \code{atof()} for \code{double} (both in \code{<stdlib.h>}). ``Read raw bytes,
    then parse'' is the same pattern you'll use for serial input later.
\end{itemize}

\subsection{Control flow}
\label{c1:a:10}
C's control-flow syntax is close to every C-family language you already know, with a few local
quirks:
\begin{itemize}
  \item Conditions must be parenthesized: \code{if (x < 10)}.
  \item Braces are technically optional for a single statement, but this codebase always uses them,
    even for one-liners: it avoids the classic bug where a second statement is added later without
    adding braces to match, silently falling outside the \code{if}:
\end{itemize}

\begin{ccode}
if (10 > x) { x = 25; }
else { printf("Out of range!\n"); }
\end{ccode}

\begin{itemize}
  \item \code{while}, \code{for}, and \code{switch} behave as expected. \code{for (;;)} or
    \code{while (true)} is the idiomatic infinite loop.
  \item \code{switch} \textbf{falls through} unless you \code{break}, a common bug source if you're
    coming from a language that doesn't do this. Give each \code{case} its own
    \code{return}/\code{break}.
  \item \code{&&}, \code{||}, \code{!} for logic; \code{==} to compare, \code{=} to assign; mixing
    them up compiles and is a classic bug (most compilers will warn on it).
  \item \textbf{Yoda notation}: writing the constant on the left of a comparison, \code{if (0 == x)}
    rather than \code{if (x == 0)}, is a defensive habit against exactly that \code{==}/\code{=}
    typo. Drop the second \code{=} with the constant on the left and \code{if (0 = x)} doesn't
    compile at all, you can't assign to a literal, so the compiler catches the mistake immediately
    instead of silently accepting an assignment that evaluates to \code{x}'s new value and is
    truthy for anything but \code{0}. This codebase writes comparisons this way throughout;
    \code{if (10 > x)} a few lines up is the same habit applied to \code{>}, even though
    inequalities don't carry the same typo risk, there's no way to accidentally assign with
    \code{>}.
\end{itemize}

\subsection{Dynamic memory allocation (\code{malloc()}/\code{realloc()}/\code{free()})}
\label{c1:a:11}
You almost certainly won't reach for this on a microcontroller: no OS to reclaim a leak, and a
fragmented heap is a liability you can't always afford. It's here because
\code{malloc()}/\code{realloc()}/\code{free()} are core C knowledge, and the double-pointer idiom
below reappears in any hosted C code (test harnesses, host-side tooling, anything running under
Linux rather than bare-metal).
\begin{itemize}
  \item \code{malloc(size)} (from \code{<stdlib.h>}) asks the heap for \code{size} bytes and returns
    a \code{void*} (a generic pointer) to it, or a null pointer if the request can't be satisfied.
    Always check the return value, and cast it to the pointer type you actually want; not required
    by C, but required by C++, and it makes the intent explicit:
\end{itemize}

\begin{ccode}
#include <stddef.h>
#include <stdlib.h>

int* data = (int*)(malloc(sizeof(int) * count));
if (NULL == data) { return 1; }
\end{ccode}

\begin{itemize}
  \item \code{realloc(ptr, new_size)} resizes an existing allocation, possibly moving it to a new
    address. On success it frees the old block for you and returns the new address; on failure it
    returns null and leaves the original block untouched. Capture the result in a temporary first,
    so a failed \code{realloc()} doesn't overwrite (and thereby leak) your only pointer to the
    still-valid original block:
\end{itemize}

\begin{ccode}
int* copy = (int*)(realloc(data, sizeof(int) * new_count));
if (NULL == copy) { return 1; } // Data is untouched and still valid on failure.
data = copy;
\end{ccode}

\begin{itemize}
  \item \code{free(ptr)} releases the block. Always set the pointer to \code{0}/\code{NULL}
    immediately after; a pointer that still holds a freed address is a \emph{dangling pointer};
    using it again is undefined behavior instead of a clean, obvious crash:
\end{itemize}

\begin{ccode}
free(data);
data = NULL;
\end{ccode}

\begin{itemize}
  \item For a function to null out the \emph{caller's} pointer, not just its own local copy of it,
    it needs the address of that pointer: the pointer-to-pointer idiom from \secref{c1:a:7}:
\end{itemize}

\begin{ccode}
/**
 * @brief Free a dynamically allocated array and null out the caller's pointer.
 *
 * @param[in,out] data Address of the pointer to the array to free.
 */
static void array_delete(int** data)
{
    if ((NULL == data) || (NULL == *data)) { return; }
    free(*data);
    *data = NULL;
}
\end{ccode}

\begin{itemize}
  \item Putting it together: \code{struct}, \code{self}, and dynamic allocation combine into the
    standard C shape for a growable array:
\end{itemize}

\begin{ccode}
#include <stddef.h>
#include <stdlib.h>

/**
 * @brief Vector structure holding floating-point numbers.
 */
typedef struct
{
    /** Data field of floating-point numbers. */
    double* data;

    /** Number of elements in the data field. */
    size_t size;
} vector_t;

/**
 * @brief Free a vector's backing storage and reset it to empty.
 *
 * @param[in,out] self Pointer to the vector to clear.
 */
static void vector_clear(vector_t* self)
{
    free(self->data);
    self->data = NULL;
    self->size = 0U;
}

/**
 * @brief Grow or shrink a vector's backing storage.
 *
 * @param[in,out] self Pointer to the vector to resize.
 * @param[in] size New number of elements.
 *
 * @return 0 on success, 1 if the allocation failed (self is left untouched).
 */
static int vector_resize(vector_t* self, const size_t size)
{
    if (NULL == self) { return 1; }

    if (0U == size)
    {
        vector_clear(self);
        return 0;
    }

    double* copy = (double*)(realloc(self->data, sizeof(double) * size));
    if (NULL == copy) { return 1; }
    self->data = copy;
    self->size = size;
    return 0;
}
\end{ccode}

\noindent\code{vector_resize()}'s \code{NULL == self} check is worth keeping, unlike a similar check
you'll see argued away later in this book (the \code{delay_ms()} of \secref{c3:a:4}, where the
pointer is only ever passed one hardcoded address from within the same file). The difference is
what kind of function this is: \code{delay_ms()} is a private, \code{static} helper with exactly
one call site the compiler can see in full; \code{vector_resize()} is written as the kind of
function a reusable module would export, called from code you don't control and can't audit by
just reading the same file. Once ``well-behaved caller'' stops being something you can verify by
inspection, checking a pointer parameter before dereferencing it stops being decorative and starts
being the thing standing between a bad call and a crash.

\subsection{Macros (\code{\#define})}
\label{c1:a:12}
You've already seen \code{#define} throughout this refresher: header guards (\secref{c1:a:6}),
\code{SENSOR_STATE_IDLE}-style constants (\secref{c1:a:3}), and every \code{LED1}/\code{F_CPU}-style
pin and clock constant this book leans on starting next chapter. It's worth understanding what it
actually does before writing your own.
\begin{itemize}
  \item \code{#define} runs in the \textbf{preprocessor}, a separate text-substitution pass that
    happens before the compiler ever sees your code. \code{#define TIMEOUT_MS 100} doesn't create a
    typed variable; it tells the preprocessor to replace every later occurrence of
    \code{TIMEOUT_MS} with the literal text \code{100}, verbatim, before compilation starts. The
    compiler never sees the name \code{TIMEOUT_MS} at all.
  \item Because it's pure text substitution, a macro has no type, no scope, and isn't visible to a
    debugger by name; it's gone by the time there's a program to debug. It also isn't namespaced to
    a file, a struct, or a function the way a variable is: once defined, \code{TIMEOUT_MS} means
    the same thing everywhere until a matching \code{#undef}, even in files that \code{#include}
    this one.
\end{itemize}

\subsubsection{Object-like vs.\ function-like macros}
\begin{itemize}
  \item An \textbf{object-like macro} stands in for a value, \code{#define F_CPU 16000000UL} or
    \code{#define LED1 1U}, the pattern this book uses constantly for pin numbers and clock
    constants.
  \item A \textbf{function-like macro} takes parameters, syntactically like a function call, but is
    still pure text substitution: no type checking on the arguments, no single-evaluation
    guarantee:
\begin{ccode}
#define SQUARE(x) (x * x)
\end{ccode}
\end{itemize}

\subsubsection{Always parenthesize: the body, and every parameter}
This is the single most important rule for writing a macro, and the reason a definition like
\code{#define TICKS_NEEDED (TIMEOUT_MS / CLOCK_TIME_MS)} is written the way it is: the whole
replacement is wrapped in parentheses. Skip that, and the macro's meaning changes depending on what
surrounds it at the call site, because substitution happens \emph{before} the usual rules of
precedence get a chance to group anything:

\begin{ccode}
#define TICKS_NEEDED TIMEOUT_MS / CLOCK_TIME_MS // No outer parentheses: fragile.

const uint16_t double_ticks = 2U * TICKS_NEEDED; // Expands to: 2U * TIMEOUT_MS / CLOCK_TIME_MS.
// * and / share precedence and group left-to-right: (2U * TIMEOUT_MS) / CLOCK_TIME_MS. Integer
// division truncates, so that equals 2U * (TIMEOUT_MS / CLOCK_TIME_MS) only if it divides evenly.
\end{ccode}

\begin{ccode}
#define TICKS_NEEDED (TIMEOUT_MS / CLOCK_TIME_MS) // Parenthesized: safe regardless of context.

const uint16_t double_ticks = 2 * TICKS_NEEDED;
// Expands to: 2 * (TIMEOUT_MS / CLOCK_TIME_MS), unambiguous no matter what precedence surrounds
// the macro at its call site.
\end{ccode}

\noindent The \code{/}-only example above only survives without the parentheses when
\code{TIMEOUT_MS} happens to be a multiple of \code{CLOCK_TIME_MS}, and that's luck, not safety;
swap in a macro built from \code{+} or \code{-} and it breaks immediately:

\begin{ccode}
#define TIMEOUT_PLUS_MARGIN TIMEOUT_MS + 10U   // No parentheses: broken.

const uint16_t scaled = 2U * TIMEOUT_PLUS_MARGIN;
// Expands to: 2U * TIMEOUT_MS + 10U, not 2U * (TIMEOUT_MS + 10U). Silently wrong, no compiler
// warning.
\end{ccode}

\noindent The same rule applies to every parameter of a function-like macro, not just the overall
body; each individual use of a parameter needs its own parentheses too:

\begin{ccode}
#define SQUARE(x) (x * x)     // Unsafe: parameter isn't parenthesized.
#define SQUARE(x) ((x) * (x)) // Safe: parameter is parenthesized everywhere it's used.

SQUARE(a + b)
// Unsafe version expands to: (a + b * a + b), wrong.
// Safe version expands to:   ((a + b) * (a + b)), correct.
\end{ccode}

\subsubsection{Macro parameters can be evaluated more than once}
A function-like macro's parameter isn't evaluated once and reused the way a real function parameter
is; the text you pass in gets substituted at every point the parameter name appears in the macro's
body. Anything with a side effect duplicates that side effect:

\begin{ccode}
#define SQUARE(x) ((x) * (x))

uint8_t i = 5U;
const uint8_t result = SQUARE(i++);   // Expands to ((i++) * (i++)): i is modified twice with
                                       // nothing to order the two, which is undefined behavior.
                                       // A real function would increment i exactly once.
\end{ccode}

\noindent Avoid passing an expression with side effects (\code{++}, \code{--}, assignment, a
function call with side effects) into a function-like macro.

\subsubsection{Multi-statement macros: wrap in \code{do \{ ... \} while (0)}}
A macro that expands to more than one statement breaks if it's ever used as the body of an
\code{if} without braces, because only the first statement ends up inside the \code{if}:

\begin{ccode}
#define LED_RESET LED_ON; LED_OFF;   // Two statements, no wrapper.

if (error) LED_RESET;   // Expands to: if (error) LED_ON; LED_OFF;
                        // LED_OFF now runs unconditionally, every time, regardless of `error`.
\end{ccode}

\noindent Wrapping the body in \code{do { ... } while (0)} turns it back into a single statement
syntactically, so it behaves correctly wherever a single statement is expected, including inside
an \code{if} with no braces:

\begin{ccode}
#define LED_RESET do { LED_ON; LED_OFF; } while (0)

if (error) LED_RESET; // Expands to a single do/while statement; behaves as one block.
\end{ccode}

\noindent This codebase always braces its \code{if}s (\secref{c1:a:10}), which sidesteps the
specific failure shown above, but the \code{do { ... } while (0)} idiom is still worth knowing:
you'll see it in other people's code, and it's the standard defense once you can't guarantee every
call site remembers to brace its \code{if}.

\subsubsection{Prefer \code{const}/\code{enum} over \code{\#define} where you can}
A macro has no type and isn't scoped, so where a real alternative exists, it's usually the better
choice:
\begin{itemize}
  \item A \code{const} variable is typed, scoped like any other variable, and visible to a debugger
    by name.
  \item An \code{enum} (\secref{c1:a:3}) groups related named integer constants the same way, with
    the same debugger visibility, and lets some compilers warn on an unhandled \code{switch} case.
\end{itemize}
Embedded C keeps reaching for \code{#define} anyway, for one specific reason: a \code{const}
variable in C is not a \emph{compile-time constant expression} the way it is in C++. Even though
its value can't change at runtime, the compiler doesn't treat it as a literal for contexts that
require one, such as a \code{case} label, the size of a file-scope array, or a \code{#if}
preprocessor check. A \code{#define} is a true compile-time constant and works in all three. An
\code{enum} constant works in the first two, but not in \code{#if}: the preprocessor runs before
the compiler and knows nothing about enums. It also pays off in generated code for something like
\code{PORTB |= (1U << LED1)}: with \code{LED1} a compile-time constant, the compiler can usually
fold the whole expression into a single bit-set instruction; with \code{LED1} a
\code{const uint8_t} variable, nothing stops it from compiling, but the shift amount is only known
at runtime, so the compiler may emit an actual runtime shift instead of folding it away. That's the
real reason this book's pin numbers and bit positions are macros instead of \code{const uint8_t}s,
not just convention.

\subsection{Bringing it together}
\label{c1:a:13}
The pattern you'll see throughout this book, and one worth internalizing before we start writing to
hardware registers next chapter, is small \code{static} functions, \code{const} everywhere
possible, explicit types, no hidden allocations, and a Doxygen comment above every public function
that isn't trivially self-explanatory:

\begin{ccode}
#include <stdio.h>


/**
 * @brief Compute resistance from voltage and current (Ohm's law).
 *
 * @param[in] voltage Voltage in V.
 * @param[in] current Current in mA.
 *
 * @return Resistance in kOhm.
 */
static inline double get_resistance(const double voltage, const double current)
{
    return voltage / current;
}

/**
 * @brief Print voltage, current, and resistance to the terminal.
 *
 * @param[in] voltage Voltage in V.
 * @param[in] current Current in mA.
 * @param[in] resistance Resistance in kOhm.
 */
static inline void print_quantities(const double voltage, const double current,
                                    const double resistance)
{
    printf("Voltage: %g V\n", voltage);
    printf("Current: %g mA\n", current);
    printf("Resistance: %g kOhm\n", resistance);
}

/**
 * @brief Application entry point.
 *
 * @return 0 on successful termination of the program.
 */
int main(void)
{
    const double voltage = 5.0;
    const double current = 0.5;
    const double resistance = get_resistance(voltage, current);
    print_quantities(voltage, current, resistance);
    return 0;
}
\end{ccode}

\noindent That style carries directly into embedded C, where you don't have a heap to hide behind
anyway.

Practice exercises for this section are in \secref{c1:sec:exercises}.
